\section{Introduction}\label{sec:intro}

\note{Arc: (1) knowledge bases are infrastructure; (2) they are written once and rarely re-read;
(3) errors are systematic because most content is produced by pipelines, and errors propagate
(IBA, orthology, ML training sets); (4) previous audits were small or statistical (RAR, CAFA,
Schnoes 2009); (5) LLM agents can now read papers and reason about ontologies, but naive LLM
curation hallucinates; (6) we show a constrained agentic system that scales, is auditable, and
finds errors that matter.}

\paragraph{Knowledge bases as infrastructure.} GO \citep{todo-go2023} \todo{scale: terms,
annotations, fraction IEA in UniProt-GOA}. Used for enrichment, CAFA benchmarks, and as labels for
protein language models \todo{cite}.

\paragraph{Write once, read rarely.} An annotation is a claim tied to a paper or pipeline at a
date. Ontology changes, retractions, and new evidence do not trigger re-review. Prior work on
misannotation: Schnoes et al.\ 2009 (enzyme superfamilies) \todo{cite}; taxon-specific rate
ratios (RAR; \citealp{todo-wei2020}); GO's own quality-control rules \todo{cite}.

\paragraph{Why errors are systematic.} Most annotations are produced by a small number of
mappings (InterPro2GO, UniProt keywords, ARBA, Ensembl Compara, PAINT/IBA). A single bad mapping
or node placement is copied to every protein that matches it. Errors are therefore correlated, and
correlated errors are invisible to evaluations that treat existing annotations as ground truth.

\paragraph{What is new here.} \todo{contributions list:}
\begin{enumerate}[nosep]
  \item A constrained agentic review architecture (schema, cached sources, verbatim quotes,
        policy validators, append-only provenance) that makes LLM curation auditable.
  \item A corpus of \nReviews{} reviews and \nAnn{} adjudicated annotations, released openly.
  \item Corpus-wide error rates stratified by evidence type, pipeline, aspect and taxon.
  \item Twelve deep dives into systematic error classes, with upstream fixes.
\end{enumerate}
